\documentclass[12pt]{article}
\usepackage{mathtools}
\usepackage{amssymb}
\usepackage{fontspec}
\setmainfont{EB Garamond}

% Set page size and margins
\usepackage[a4paper,top=2cm,bottom=1.5cm,left=1cm,right=1cm]{geometry}
\usepackage{fancyhdr}
\pagestyle{fancy}
\setlength{\parindent}{0pt}

% Useful packages
\usepackage{minted}
\usepackage{graphicx}
\graphicspath{{./media/}}
\usepackage{subfig}
\usepackage{float}
\usepackage{siunitx}
\usepackage[hidelinks]{hyperref}
\usepackage{multirow}

% Change "Table" to "Πίνακας"
\renewcommand{\tablename}{Πίνακας}


\title{\vspace{-2cm} Θεμελιώσεις Κρυπτογραφίας \\
        Homework 4}
\author{Μαρινόπουλος Χριστόφορος 10522}
\date{Εαρινό εξάμηνο 2025}

\lhead{Homework 4}
\chead{Θεμελιώσεις Κρυπτογραφίας}
\rhead{Μαρινόπουλος Χριστόφορος}
\begin{document}
\maketitle
\section*{Ερώτημα 1}
Έστω οι συναρτήσεις Bool, \( f, g : \mathbb{F}_2^n \rightarrow \{0,1\} \)

Ονομάζουμε ``correlation coefficient''

\[
C(f, g) = \frac{1}{2^n} \sum_{\mathbf{x} \in \mathbb{F}_2^n} (-1)^{f(\mathbf{x}) + g(\mathbf{x})}
\]

Αν, \quad 
\( f, g : \mathbb{F}_2^3 \rightarrow \{0,1\} \), \quad 
\( f = x_1 x_2 + x_3 \), \quad 
\( g = x_1 + x_3 \)

Υπολογίστε τον \( C(f, g) \).

\section*{Απάντηση}
Για τις συγκεκριμένες συναρτήσεις ισχύει ότι:
\[
C_i = \frac{1}{2^3} (-1)^{f(x) + g(x)} =
\begin{cases}
\frac{1}{8}, & \text{αν } f(x) = g(x) \\
-\frac{1}{8}, & \text{αν } f(x) \ne g(x)
\end{cases}
\]




Στον \textbf{\protect\hyperref[πίνακας 1]{Πίνακα 1}} παρουσιάζονται οι τιμές των \( f(x) \) και \(g(x)\) υπολογισμένες για κάθε \( x_1 , x_2 , x_3 \) καθώς και ο \(C_i\).

\renewcommand{\arraystretch}{1.8} % αύξηση ύψους γραμμών

\begin{table}[h!]
\centering
\caption{Τιμές \( f(x) \), \( g(x) \) και \( C_i = \frac{1}{8}(-1)^{f(x) + g(x)} \)}
\label{πίνακας 1}
\begin{tabular}{|c||c|c|c||c|c||c|}
\hline
\( i \) & \( x_1 \) & \( x_2 \) & \( x_3 \) & \( f(x) \) & \( g(x) \) & \( C_i \) \\
\hline
1 & 0 & 0 & 0 & 0 & 0 & \( \frac{1}{8} \) \\
\hline
2 & 0 & 0 & 1 & 1 & 1 & \( \frac{1}{8} \) \\
\hline
3 & 0 & 1 & 0 & 0 & 0 & \( \frac{1}{8} \) \\
\hline
4 & 0 & 1 & 1 & 1 & 1 & \( \frac{1}{8} \) \\
\hline
5 & 1 & 0 & 0 & 0 & 1 & \( -\frac{1}{8} \) \\
\hline
6 & 1 & 0 & 1 & 1 & 1 & \( \frac{1}{8} \) \\
\hline
7 & 1 & 1 & 0 & 1 & 1 & \( \frac{1}{8} \) \\
\hline
8 & 1 & 1 & 1 & 1 & 1 & \( \frac{1}{8} \) \\
\hline
\end{tabular}
\end{table}
 
Για να υπολογίσουμε τον \( C(f,g) \) αρκεί να υπολογίσω το:\( \sum_{i=1}^{8}C_i\)

Τελικά:
\[
C(f,g) = \sum_{i=1}^{8}C_i = \mathbf{\frac{6}{8}}
\]
\hfill \blacksquare

\section*{Ερώτημα 2}
Έστω \( S \) ένα \( S \) κιβώτιο, \( S : \mathbb{F}_2^n \rightarrow \mathbb{F}_2^m \)

Ονομάζουμε ``Differential Uniformity''

\[
\mathrm{Diff}(S) = \max_{\substack{\mathbf{x} \in \mathbb{F}_2^n - \{0\} \\ \mathbf{y} \in \mathbb{F}_2^m}} 
\left| \left\{ \mathbf{z} \in \mathbb{F}_2^n : S(\mathbf{z} \oplus \mathbf{x}) \oplus S(\mathbf{z}) = \mathbf{y} \right\} \right|
\]

Αν, το \( S \) είναι το \( S \) -Box του \textbf{\protect\hyperref[πίνακας 2]{  Πίνακα 2}}, βρείτε τον \( \mathrm{Diff}(S) \).

\begin{table}[h!]
\label{πίνακας 2}
\caption{S-box}
\centering
\resizebox{\textwidth}{!}{%
\begin{tabular}{|c|c|cccccccccccccccc|}
\hline
\multicolumn{2}{|c|}{\textbf{S₅}} & \multicolumn{16}{c|}{\textbf{Middle 4 bits of input}} \\
\multicolumn{2}{|c|}{} & 0000 & 0001 & 0010 & 0011 & 0100 & 0101 & 0110 & 0111 & 1000 & 1001 & 1010 & 1011 & 1100 & 1101 & 1110 & 1111 \\
\hline
\multirow{4}{*}{\rotatebox{90}{\textbf{Outer bits}}} 
& 00 & 0010 & 1100 & 0100 & 0001 & 0111 & 1010 & 1011 & 0110 & 1000 & 0101 & 0011 & 1111 & 1101 & 0000 & 1110 & 1001 \\
& 01 & 1110 & 1011 & 0010 & 1100 & 0100 & 1111 & 0111 & 0001 & 0101 & 1000 & 1101 & 1010 & 1001 & 1001 & 0000 & 0110 \\
& 10 & 0100 & 1100 & 0001 & 0111 & 1010 & 1110 & 0011 & 0100 & 0111 & 1000 & 1111 & 0000 & 1101 & 0010 & 0000 & 1110 \\
& 11 & 1011 & 1000 & 0000 & 0111 & 0111 & 1110 & 0010 & 1101 & 1111 & 0000 & 1011 & 1010 & 0100 & 0101 & 0101 & 0011 \\
\hline
\end{tabular}%
}
\end{table}

\section*{Απάντηση}

\section*{Απάντηση}

Για τον υπολογισμό του \( \mathrm{Diff}(S) \) για το S-Box του 
\textbf{\protect\hyperref[πίνακας 2]{Πίνακα 2}}, αρκεί να κατασκευαστεί ο διαφορικός 
πίνακας κατανομής του. Από αυτόν, επιλέγεται η μέγιστη τιμή (εκτός της περίπτωσης με 
μηδενική διαφορά εισόδου), η οποία αντιστοιχεί στην τιμή του Differential Uniformity.
\newline
\newline
Η κατασκευή του πίνακα απαιτεί την εξέταση όλων των δυνατών διαφορών εισόδου \( \Delta x 
\neq 0 \) και τον υπολογισμό των αντίστοιχων διαφορών εξόδου \( \Delta y \), για κάθε τιμή 
εισόδου. Ο συνολικός αριθμός πράξεων που απαιτούνται είναι: \( (64 - 1) \times 64 \times 3 
= 12096 \). (Στην πράξη, ο αριθμός αυτός είναι μικρότερος, εφόσον αξιοποιηθούν συμμετρικές ιδιότητες ή επαναλαμβανόμενα αποτελέσματα μεταξύ υπολογισμών.)
\newline
\newline

Έτσι, για την εκτέλεση των υπολογισμών, αναπτύχθηκε κωδικας ο οποίος βρίσκεται στο αρχείο \texttt{homework\_4\_10522.py}.
\newline
\newline
Με την εκτέλεση του, είναι προφανές: \( \mathbf{\mathrm{Diff}(S) = 22} \).


\hfill \blacksquare

\end{document}